\originalpreface
粒子数可以改变, 同时局域操作要遵守相对论因果性: 这两个要求把量子力学引向量子场论. 本册读者已会态空间、谐振子和基本微扰论, 但不预设学过场论. 第1--2章补入四维记号、局域场、作用量与 Noether 定理, 第3--12章建立自由场到 QED 的计算链, 第13--19章进入非 Abel 规范场、谱与 LSZ、重整化和尺度流. 后续作业与理论合订, 不是原英文 PDF 的替代目录.

主来源为 MIT OpenCourseWare 的 Hong Liu 两门研究生课程: \href{https://ocw.mit.edu/courses/8-323-relativistic-quantum-field-theory-i-spring-2023/}{8.323 Relativistic Quantum Field Theory I, Spring 2023}, 以及 \href{https://ocw.mit.edu/courses/8-324-relativistic-quantum-field-theory-ii-fall-2010/}{8.324 Relativistic Quantum Field Theory II, Fall 2010}. 两课均每周两次90分钟讲课、一次60分钟习题课. 原课先修分别为8.321, 以及8.322和8.323; 本册只补与当前推导直接相关的最短先修, 不把仓库已有本科量子力学卷等同于完成全部研究生量子理论先修.

本轮冻结的理论范围是8.323的26讲主题和8.324实际公开的26份讲义. 考核范围为8.323十二套与8.324七套公开作业及匹配官方解. 两课原大纲均明确没有考试, 成绩来自作业; 8.323最低一次作业成绩不计, 8.324大纲称八次作业但公开页仅有1--7. 因而本册不提供虚构的期中、期末或8.324第8套. 仅有商业教材题号的项目在题旁保留引用和缺口, 不从答案反造题面. 九份8.323习题课笔记另行归档供补充参考, 第9次原件明确不可得, 并不宣称已逐份翻译全部习题课笔记或逐分钟转录视频.

全书取 $c=\hbar=1$, 度规 $(+---)$; 原课多取 $(-+++)$. 原题的数学与物理条件保留, 涉及符号的转换逐处交代. 中文理论是按概念重新组织的推导, 作业解是公开官方解的中文重编与独立补充, 确定的原件排印问题保留定位后修正. 非微扰场论的严格构造、一般自旋统计/CPT定理的完整公理证明、禁闭与完整标准模型不在本轮证明范围.

\noindent\textbf{8.323 原课进度.} 下表据原 \href{https://ocw.mit.edu/courses/8-323-relativistic-quantum-field-theory-i-spring-2023/pages/calendar/}{Calendar} 按每两讲排列, 不把讲次反推为未提供的具体日期. 课历在第2讲布置作业1, 在第4,6,\ldots,24讲依次收作业1--11并布置下一套, 第26讲收作业12.
\begin{longtable}{@{}p{.10\linewidth}p{.71\linewidth}p{.13\linewidth}@{}}
\toprule 讲次&主题&本册理论章\\\midrule\endhead
1--2&经典场、局域性; 对称性与守恒律&1--2\\
3--4&场论的必要性; 自由标量正则量子化&1,3\\
5--6&复标量与反粒子; 传播子与 Green 函数&3--4\\
7--8&相互作用与散射矩阵; 量子力学路径积分&5--6\\
9--10&场路径积分; 时间排序关联函数&5\\
11--12&微扰关联函数; Feynman 图与组合系数&6\\
13--14&Dirac 方程; Lorentz 协变性&7\\
15--16&Dirac 经典解; 正则量子化&7--8\\
17--18&手征与 Majorana 旋量; 离散对称&8--9\\
19--20&费米子路径积分; Maxwell 正则量子化&9--10\\
21--22&量子 Maxwell 场; QED&10\\
23--24&截面与衰变率; QED基本过程之一&11\\
25--26&QED基本过程之二; 量子涨落与重整化&11--12\\\bottomrule
\end{longtable}

\noindent\textbf{8.324 原课进度.} 下表据 \href{https://ocw.mit.edu/courses/8-324-relativistic-quantum-field-theory-ii-fall-2010/pages/syllabus/}{Syllabus 内 Calendar}. 当前公开资料没有逐日的作业收发日期, 因而这里只保留真实讲次顺序.
\begin{longtable}{@{}p{.10\linewidth}p{.71\linewidth}p{.13\linewidth}@{}}
\toprule 讲次&主题&本册理论章\\\midrule\endhead
1--2&对称性、Lie群与代数; 规范原理&13\\
3--4&Yang--Mills 理论及其非 Abel 结构&13\\
5--6&规范场量子化; BRST、物理态与幺正性&14\\
7--8&场与质量重整化; 标量一圈实例&15\\
9--10&紫外发散减除; 不稳定粒子与共振&15\\
11--12&散射矩阵与 LSZ; 重整化 QED&15--16\\
13--14&重整化 QED; 顶角函数&16\\
15--16&异常磁矩; 真空极化&16\\
17--18&表面发散次数; 发散抵消&17\\
19--20&一般重整化问题; Wilson 方法&17--18\\
21--22&重整化群流及精确流方程&18\\
23--24&尺度流; 传统方案的 beta 函数&18--19\\
25--26&beta函数、六维标量例、QED与QCD尺度运行&19\\\bottomrule
\end{longtable}

原课程、讲义作者 Hong Liu, 以及署名的习题课作者 Wentao Cui 的信用均保留. 本册是非官方中文改编, 不表示 MIT 审定或认可. 原件与中文改编按 \href{https://creativecommons.org/licenses/by-nc-sa/4.0/}{CC BY-NC-SA 4.0} 非商业共享, 具体来源、第三方例外核查和不可得材料见册末及源码包许可说明. 商业教材本体、字体和视频没有打包.
